% =============================================================================
% 6. Considerações Parciais e Próximos Passos
% =============================================================================

\section{Considerações Parciais e Próximos Passos}

Até esta entrega parcial foram consolidados: (i) a fundamentação teórica cobrindo modelos de espaço vetorial, \textit{embeddings} densos, busca aproximada via HNSW, arquiteturas de bancos de dados vetoriais, RAG e metodologia de \textit{benchmarking}; (ii) um ambiente experimental reproduzível com três SGBDs operando em condições controladas via Docker Compose, com versões fixadas e dependências pinadas; (iii) \textit{scripts} de \textit{benchmark} dos cenários A e B implementados e validados em TDD, com 198 testes automatizados; e (iv) resultados quantitativos dos cenários A e B em \textit{datasets} de 100 mil e 500 mil vetores, com curvas \textit{recall} $\times$ QPS obtidas para os três sistemas.

O código-fonte está disponível em repositório Git com integração contínua, e os arquivos de resultado de cada execução são versionados junto ao código, de modo que todo valor reportado neste texto pode ser rastreado até a medição que o originou.

No Cenário A, os três sistemas ultrapassam \textit{recall@10} de 0,99 em 100 mil vetores com $ef\_search \geq 128$; em 500 mil, o mesmo ponto de operação já não é suficiente para todos, e o Weaviate requer $ef\_search = 256$. Os perfis diferem de forma consistente: o Weaviate opera com as menores latências e o maior \textit{throughput}, o Qdrant apresenta a distribuição de latência mais concentrada, e o pgvector é o mais sensível ao aumento da base, com queda acentuada de \textit{throughput} sob $ef\_search$ elevado em 500 mil vetores.

O Cenário B produziu o resultado metodologicamente mais instrutivo desta fase. A leitura inicial --- \textit{recall} perfeito dos sistemas especializados sob filtros restritivos --- não se sustentou: o valor decorria da substituição automática da busca aproximada por varredura exata quando o subconjunto elegível é pequeno, e não da qualidade da filtragem. Reexecutado o cenário com essa substituição desativada, o comportamento se reorganiza: o Qdrant preserva \textit{recall} elevado, o Weaviate cai a 0,5570 sob seletividade de 1\%, e o pgvector permanece limitado pela estratégia de \textit{post-filtering}, sem que a indexação do atributo de filtro altere o quadro. A conclusão que se sustenta é mais restrita e mais informativa que a inicial: há diferença arquitetural real na busca vetorial com filtro, mas ela só é mensurável depois de neutralizado o efeito das heurísticas de execução de cada sistema.

As etapas seguintes do cronograma compreendem: (i) a execução dos experimentos em escala de 1 milhão de vetores, que permitirá observar o ponto a partir do qual cada arquitetura se diferencia de forma mais pronunciada; (ii) o cenário de carga mista (Cenário C), que simulará inserções concorrentes com buscas, característico de aplicações RAG em produção; e (iii) a análise comparativa final com sistematização de recomendações práticas por cenário de uso. A possível extensão para o \textit{Cluster} HPC do IEG/UFOPA, condicionada à viabilidade técnica e à disponibilidade de acesso, será avaliada nessa segunda metade do projeto e, caso confirmada, incorporada ao relatório final previsto para dezembro de 2026.
